\section{Limiti del primo approccio e scelta del metodo attuale}
\label{sec:limiti}

Nella versione del notebook \texttt{2\_1\_benchmark\_thor.ipynb} salvata nel commit \texttt{52cd3d8} (11 luglio, 19:03), il prompt per il LLM viene ancora costruito campionando \texttt{n\_rows=2000} righe del dataset in formato CSV grezzo --- lo stesso approccio delle fasi precedenti, ma con un campione dieci volte più grande rispetto alle 300 righe usate con o3-mini. I quattro file di risultati salvati nella stessa finestra di lavoro (tra le 18:37 e le 19:11 dello stesso giorno) mostrano un consumo di token in ingresso molto variabile (da 5131 a 31\,866 a seconda della run) e risultati poco stabili: in una delle run, ad esempio, gli errori relativi registrati raggiungono il 785\% su SE e il 110\% su IE, mentre un'altra produce SE, DE e IE tutti nulli. Il repository non consente di stabilire con precisione la causa esatta di questa variabilità (temperatura, gestione del contesto del server, o altro), ma il pattern osservabile --- risultati incoerenti tra run nominalmente equivalenti, a fronte di un prompt che nel caso di 2000 righe supera abbondantemente le 25\,000--30\,000 token --- rende evidente che inviare il dataset grezzo al modello non fosse una strategia praticabile, sia in termini di costo computazionale sia di affidabilità.

Nove minuti dopo l'ultimo di questi tentativi, il commit \texttt{b3ee1f3} (11 luglio, 19:13, \emph{``FIXED: aggregate dataset into conditional probability tables instead of raw rows''}) riscrive la costruzione del prompt: al posto delle righe campionate, il notebook calcola con \texttt{groupby} di pandas cinque tabelle di probabilità condizionale pre-aggregate (P(Y|X), P(Z), P(Y|X,Z), P(W|X,Z), P(Y|X,W,Z)) e le inserisce nel prompt insieme alle formule di identificazione. Il file di risultati prodotto nello stesso commit (\texttt{adult\_20260711\_191135.json}) mostra un netto miglioramento di stabilità: input ridotto a 6748 token (contro le decine di migliaia precedenti) e un tempo di risposta di circa 4 secondi. Questo approccio --- tabelle pre-aggregate anziché dati grezzi --- diventa da questo momento il metodo standard adottato in tutti gli esperimenti successivi, incluso il confronto finale descritto in \S\ref{sec:confronto-finale}.

È opportuno notare che questo cambio di strategia nella costruzione del prompt avviene in concomitanza con il passaggio, già introdotto in \S\ref{sec:thor}, da un modello cloud (o3-mini via API OpenAI) a un modello locale open-weight (Qwen2.5, servito via \texttt{llama.cpp} su Thor): il repository documenta con certezza il cambio di formato del prompt (dati grezzi $\to$ tabelle aggregate) come reazione diretta ai problemi di costo e stabilità osservati; la relazione causale con la scelta del modello locale, motivata presumibilmente anche da considerazioni di costo e disponibilità di calcolo GPU dedicato, non è invece esplicitata nei messaggi di commit e viene qui presentata come inferenza plausibile, non come fatto direttamente documentato.
